\documentclass[11pt]{exam}
\usepackage{foerderkurs}

\begin{document}

\werkzeug{W1}{Sprache und Mathematik}

Viele Schwierigkeiten in Mathematik sind eigentlich Schwierigkeiten mit der
Sprache: Ein Wort bedeutet in der Mathematik etwas anderes als im Alltag, ein
Satz wird Wort für Wort übersetzt statt dem Sinn nach, oder ein Ergebnis wird
nicht zurück in die Situation übertragen. Diese Einheit besteht aus drei
Teilen, die über mehrere Wochen verteilt behandelt werden.

\begin{lehrperson}[Überblick]
\textbf{Aufbau:} drei Teile zu je ca.\ 15 Minuten Instruktion, danach Übungen
in der Übungszeit. Die Teile werden über das Jahr verteilt, z.\,B.:

\begin{tabular}{@{}llp{7.2cm}@{}}
  Teil 1 & Wörter und Operatoren   & früh im Algebra-Fundament, zusammen mit Mini-Lektion 0.1 \\
  Teil 2 & Vom Text zur Gleichung  & vor den linearen Gleichungen (Mini-Lektionen 6.1/6.2) \\
  Teil 3 & Von der Mathematik zurück zum Text & in der Phase Proportionalität und Funktionen
\end{tabular}

\textbf{Wichtiger als die Teile selbst:} Die Routine aus Teil 2 und die
Antwortsätze aus Teil 3 werden danach bei \emph{jeder} Textaufgabe im Kurs
wieder eingefordert. Erst die Wiederholung macht daraus ein Werkzeug.

\textbf{Verknüpfung:} Standortbestimmung Strang 9 (Aufgaben 9.1--9.6) sowie
4.2, 7.3, 7.4. Hinweis zur Diagnose: Wer eine Aufgabe als reine Rechnung löst,
dieselbe Aufgabe als Text aber nicht, hat ein Sprachproblem und kein
Rechenproblem.

\textbf{Deutsch als Zweitsprache:} Für Schülerinnen und Schüler mit einer
anderen Erstsprache lohnt es sich, im Glossar eine dritte Spalte für die
Übersetzung in die Erstsprache anzubieten.
\end{lehrperson}

% ===========================================================================
\werkzeugteil{1}{Wörter und Operatoren}
  {die wichtigsten mathematischen Fachwörter verstehen und weiss, was eine
   Aufgabe mit «berechnen», «begründen» oder «zeigen» von mir verlangt.}

\begin{einstieg}
Viele Wörter haben im Alltag eine andere Bedeutung als in der Mathematik.
Ergänzen Sie die Tabelle.

\begin{center}
\renewcommand{\arraystretch}{1.7}
\newcommand{\lsg}[1]{\ifprintanswers\small #1\fi}
\begin{tabular}{|l|p{5.3cm}|p{5.3cm}|}
  \hline
  \textbf{Wort} & \textbf{im Alltag} & \textbf{in der Mathematik} \\ \hline
  Produkt   & \lsg{Ware, Erzeugnis}         & \lsg{Ergebnis einer Multiplikation} \\ \hline
  Differenz & \lsg{Unterschied, Streit}     & \lsg{Ergebnis einer Subtraktion} \\ \hline
  Lösung    & \lsg{Ausweg; Flüssigkeit}     & \lsg{Zahl, die eine Gleichung erfüllt} \\ \hline
  Mittel    & \lsg{Hilfsmittel, Medikament} & \lsg{Durchschnitt} \\ \hline
  Bruch     & \lsg{Knochenbruch, Regelbruch}& \lsg{Zahl wie $\frac{3}{4}$} \\ \hline
\end{tabular}
\end{center}
\end{einstieg}

\begin{lehrperson}[Teil 1, Einstieg]
Zu zweit 3 Minuten lang ausfüllen lassen, dann sammeln. Ziel ist nicht
Vollständigkeit, sondern die Einsicht: Mathematik ist eine eigene Sprache mit
eigenen Wörtern. Schülerinnen und Schüler können weitere Beispiele nennen
(z.\,B.\ Wurzel, Potenz, Umfang, Term).
\end{lehrperson}

\begin{merke}
\textbf{Die vier Grundbegriffe.} Bei Differenz und Quotient kommt es auf die
Reihenfolge an.

\begin{center}
\renewcommand{\arraystretch}{1.3}
\begin{tabular}{llll}
  \textbf{Summe}      & $a + b$           & «die Summe von $a$ und $b$»      & $a$, $b$: Summanden \\
  \textbf{Differenz}  & $a - b$           & «die Differenz von $a$ und $b$»  & \\
  \textbf{Produkt}    & $a \cdot b$       & «das Produkt von $a$ und $b$»    & $a$, $b$: Faktoren \\
  \textbf{Quotient}   & $a \geteilt b$    & «der Quotient von $a$ und $b$»   & $a$: Dividend, $b$: Divisor
\end{tabular}
\end{center}

\textbf{Vergleichen.} Achten Sie auf die Reihenfolge bei «weniger als»:

\begin{center}
\renewcommand{\arraystretch}{1.3}
\begin{tabular}{ll@{\qquad}ll}
  um $3$ grösser als $x$      & $x + 3$  & $3$-mal so gross wie $x$ & $3 \cdot x$ \\
  um $3$ kleiner als $x$      & $x - 3$  & die Hälfte von $x$       & $\frac{x}{2}$ \\
  $3$ weniger als $x$         & $x - 3$  & höchstens $3$            & $\leq 3$ \\
  $x$ vermindert um $3$       & $x - 3$  & mindestens $3$           & $\geq 3$
\end{tabular}
\end{center}
\end{merke}

\begin{merke}
\textbf{Operatoren} sagen, was eine Aufgabe von Ihnen erwartet:

\begin{center}
\renewcommand{\arraystretch}{1.3}
\begin{tabular}{@{}lp{11cm}@{}}
  \textbf{angeben, nennen}   & nur das Ergebnis, ohne Begründung \\
  \textbf{berechnen}         & Ergebnis \emph{mit} nachvollziehbarem Rechenweg \\
  \textbf{bestimmen}         & Ergebnis finden; der Weg (Rechnung, Zeichnung) ist frei wählbar \\
  \textbf{vereinfachen}      & einen Term so umformen, dass er kürzer oder übersichtlicher wird \\
  \textbf{begründen}         & in Worten erklären, \emph{warum} etwas gilt \\
  \textbf{zeigen, nachweisen}& das Ergebnis ist bekannt; der Weg dorthin ist gefragt \\
  \textbf{skizzieren}        & ungefähr, die wichtigsten Merkmale, ohne Massstab \\
  \textbf{zeichnen}          & genau, mit Massstab und Hilfsmitteln
\end{tabular}
\end{center}
\end{merke}

\begin{lehrperson}[Operatoren]
Die Operatoren werden nicht an allen Schulen und von allen Lehrpersonen gleich
streng verwendet. Es lohnt sich, die Liste mit den Kolleginnen und Kollegen
abzugleichen. Für die Schülerinnen und Schüler ist vor allem der Unterschied
zwischen \emph{angeben}, \emph{berechnen} und \emph{begründen} wichtig, weil
er an Prüfungen über Punkte entscheidet.
\end{lehrperson}

\begin{schreibweisen}
\begin{itemize}
  \item «Das Produkt \emph{von} $3$ und $5$» und «das Produkt \emph{aus} $3$
        und $5$» bedeuten dasselbe.
  \item In englischen Videos: sum, difference, product, quotient; «less
        than» hat dieselbe Falle: «$3$ less than $x$» ist $x - 3$.
  \item KI-Programme verwenden die Operatoren oft weniger genau. Ein
        «Beweis» in einer KI-Antwort ist manchmal nur ein Beispiel.
\end{itemize}
\end{schreibweisen}

\begin{questions}

\question Schreiben Sie als Rechnung und berechnen Sie das Ergebnis.
\begin{parts}
  \part der Quotient von $48$ und $6$
  \part die Differenz von $7$ und $12$
  \part das Produkt aus der Summe von $3$ und $5$ und der Zahl $4$
  \part die Summe aus dem Produkt von $3$ und $5$ und der Zahl $4$
\end{parts}
\begin{solution}
  a) $48 \geteilt 6 = 8$ \quad b) $7 - 12 = -5$ \quad
  c) $(3 + 5) \cdot 4 = 32$ \quad d) $3 \cdot 5 + 4 = 19$
\end{solution}
\raum{3cm}

\question Schreiben Sie als Term.
\begin{multicols}{2}
\begin{parts}
  \part um $7$ grösser als $x$
  \part $7$-mal so gross wie $x$
  \part $7$ weniger als $x$
  \part $x$ vermindert um das Doppelte von $y$
  \part das Quadrat der Summe von $a$ und $b$
  \part die Summe der Quadrate von $a$ und $b$
\end{parts}
\end{multicols}
\begin{solution}
  a) $x + 7$ \quad b) $7x$ \quad c) $x - 7$ \quad d) $x - 2y$ \quad
  e) $(a + b)^2$ \quad f) $a^2 + b^2$
\end{solution}
\raum{3cm}

\question Was erwartet die Aufgabe von Ihnen? Beschreiben Sie in einem Satz,
was Sie aufschreiben müssen. Sie müssen die Aufgaben \emph{nicht} lösen.
\begin{parts}
  \part Geben Sie die Lösung der Gleichung $2x = 10$ an.
  \part Berechnen Sie den Flächeninhalt eines Rechtecks mit den Seiten
        $4$\,cm und $7$\,cm.
  \part Begründen Sie, warum $2a + 3b$ nicht gleich $5ab$ ist.
  \part Zeigen Sie, dass $(a + b)^2 = a^2 + 2ab + b^2$ gilt.
\end{parts}
\begin{solution}
  a) Nur das Ergebnis: $x = 5$. \
  b) Rechnung und Ergebnis: $4 \cdot 7 = 28$, also $28$\,cm$^2$. \
  c) Eine Erklärung in Worten, z.\,B.\ mit einem Gegenbeispiel: Für $a = 1$
  und $b = 2$ ist $2a + 3b = 8$, aber $5ab = 10$. (Achtung: $a = b = 1$ ergibt
  zufällig zweimal $5$ und ist darum kein Gegenbeispiel.) \
  d) Eine Herleitung, z.\,B.\ durch Ausmultiplizieren von $(a + b) \cdot (a + b)$.
\end{solution}
\raum{3.5cm}

\end{questions}

\begin{lehrperson}[Übung 3c]
Bei 3c wählen viele zuerst $a = b = 1$ und erhalten zufällig zweimal $5$. Das ist ein guter Anlass zu besprechen, dass ein
einzelnes passendes Beispiel nichts beweist, ein einzelnes Gegenbeispiel aber
genügt, um eine Aussage zu widerlegen.
\end{lehrperson}

\begin{transfer}
Operatoren gibt es in allen Fächern: In Geschichte heisst «erläutern» etwas
anderes als «nennen», in Biologie «beschreiben» etwas anderes als
«erklären». Wer an einer Prüfung zuerst den Operator markiert, weiss, wie
ausführlich die Antwort sein muss.
\end{transfer}

\begin{tcolorbox}[fkbasis, colframe=fkgrau, colback=white, title={Ihr Glossar}]
Legen Sie auf den letzten Seiten Ihres Hefts ein \textbf{Glossar} an: Wort,
Bedeutung in der Mathematik, ein Beispiel. Ergänzen Sie es jedes Mal, wenn
Ihnen im Kurs ein neues Fachwort begegnet.
\end{tcolorbox}

\begin{exitfrage}
Schreiben Sie als Term: «$5$ weniger als das Dreifache von $y$».
\ifprintanswers\par\medskip\textbf{Lösung:} $3y - 5$
\else\par\vspace{1.5cm}\fi
\end{exitfrage}

% ===========================================================================
\newpage
\werkzeugteil{2}{Vom Text zur Gleichung}
  {eine Textaufgabe mit einer festen Routine in eine Gleichung übersetzen und
   meine Gleichung mit einer Zahlenprobe überprüfen.}

\begin{einstieg}
An einer Schule gibt es sechsmal so viele Schülerinnen und Schüler wie
Lehrpersonen. Schreiben Sie diese Aussage als Gleichung. Verwenden Sie $S$
für die Anzahl Schülerinnen und Schüler und $L$ für die Anzahl Lehrpersonen.
\ifprintanswers\else\par\vspace{1.5cm}\fi
\end{einstieg}

\begin{lehrperson}[Teil 2, Einstieg]
Diese Aufgabe stammt aus einer bekannten Studie mit Studierenden; auch dort
schrieb ein grosser Teil $6S = L$. Der Fehler entsteht durch
Wort-für-Wort-Übersetzung: «sechsmal $S$ gleich $L$». Antworten sammeln,
ohne aufzulösen, und dann die Zahlenprobe aus dem Merkkasten vorführen:
Bei $L = 10$ Lehrpersonen müsste es $60$ Schülerinnen und Schüler geben. Setzt
man das in $6S = L$ ein, erhält man $360 = 10$. Die Gleichung beschreibt also
nicht die Situation.
\end{lehrperson}

\begin{merke}
\textbf{Routine in vier Schritten}
\begin{enumerate}
  \item \textbf{Lesen:} zweimal, einmal für den Überblick und einmal genau.
  \item \textbf{Markieren:} Was ist gegeben? Was ist gesucht? Was ist
        überflüssig?
  \item \textbf{Übersetzen:} Zuerst die Variable festlegen, und zwar als
        \emph{Zahl}: «$x$: Anzahl Lehrpersonen», nicht «$x$ = Lehrpersonen».
        Dann Satz für Satz in Mathematik übersetzen.
  \item \textbf{Prüfen:} Zahlenprobe mit einfachen Zahlen, Gleichung zurück in
        Worte übersetzen, Ergebnis auf Plausibilität prüfen.
\end{enumerate}
\end{merke}

\begin{beispiel}[1: Die Zahlenprobe]
«Es gibt sechsmal so viele Schülerinnen und Schüler wie Lehrpersonen.»

Zahlenprobe: Bei $10$ Lehrpersonen gibt es $60$ Schülerinnen und Schüler.
Die Anzahl Schülerinnen und Schüler ist also das Sechsfache der Anzahl
Lehrpersonen:
\[ S = 6 \cdot L \qquad\text{Probe: } 60 = 6 \cdot 10 \;\checkmark \]
\falsch\quad $6 \cdot S = L$ \qquad Probe: $6 \cdot 60 = 360 \neq 10$
\end{beispiel}

\begin{beispiel}[2: Die ganze Routine]
«Anna ist $4$ Jahre älter als Ben. Zusammen sind sie $30$ Jahre alt. Wie alt
ist Anna?»
\begin{enumerate}
  \item[1.] Lesen: zweimal.
  \item[2.] Gegeben: Altersunterschied $4$, Summe $30$. Gesucht: Annas Alter.
  \item[3.] $x$: Alter von Ben in Jahren. Dann ist Anna $x + 4$ Jahre alt.
            \[ x + (x + 4) = 30 \quad\Rightarrow\quad 2x + 4 = 30
               \quad\Rightarrow\quad x = 13 \]
            Anna ist $13 + 4 = 17$ Jahre alt.
  \item[4.] Probe: $17 + 13 = 30$ \checkmark, und $17$ ist um $4$ grösser als
            $13$ \checkmark.
\end{enumerate}
\end{beispiel}

\begin{beispiel}[3: Schlüsselwörter können täuschen]
«Nach einem Rabatt von $20\,\%$ kostet ein Velo insgesamt CHF 600. Wie viel
kostete es vorher?»

Das Wort «insgesamt» verleitet zum Addieren, und «Rabatt» zum Abziehen von
$20\,\%$ von CHF 600. Beides ist falsch. Übersetzen statt Schlüsselwörter
suchen:

$x$: Preis vorher in CHF. Nach dem Rabatt bleiben $80\,\%$ davon:
\[ 0.8 \cdot x = 600 \quad\Rightarrow\quad x = 750 \]
Probe: $20\,\%$ von CHF 750 sind CHF 150, und $750 - 150 = 600$ \checkmark
\end{beispiel}

\begin{lehrperson}[Beispiele]
Beispiel 3 knüpft an Aufgabe 9.2\,b der Standortbestimmung an (Schlüsselwort
«mal»). Viele Schülerinnen und Schüler haben in der Primarschule gelernt,
nach Signalwörtern zu suchen («insgesamt» heisst plus, «weniger» heisst
minus). Diese Strategie funktioniert bei einfachen Aufgaben und versagt bei
allen anderen. Es lohnt sich, das offen anzusprechen.
\end{lehrperson}

\begin{questions}

\question Schreiben Sie als Gleichung. Legen Sie zuerst die Variablen fest und
machen Sie danach die Zahlenprobe.
\begin{parts}
  \part In einem Parkhaus stehen dreimal so viele Autos wie Motorräder.
  \part Lea hat CHF 5 mehr als Tim.
  \part Ein Kilo Äpfel kostet halb so viel wie ein Kilo Kirschen.
  \part Auf $4$ Katzen kommen $7$ Hunde.
\end{parts}
\begin{solution}
  a) $A$: Anzahl Autos, $M$: Anzahl Motorräder; $A = 3 \cdot M$ \
  b) $L$, $T$: Betrag in CHF; $L = T + 5$ \
  c) $A$, $K$: Preis pro Kilo in CHF; $A = \frac{K}{2}$ oder $2 \cdot A = K$ \
  d) $K$: Anzahl Katzen, $H$: Anzahl Hunde; $7 \cdot K = 4 \cdot H$
  (Probe: $K = 4$, $H = 7$ ergibt $28 = 28$) oder $\frac{K}{H} = \frac{4}{7}$
\end{solution}
\raum{4cm}

\question Lösen Sie mit der Routine. Schreiben Sie am Schluss einen
Antwortsatz.
\begin{parts}
  \part Subtrahiert man $5$ vom Dreifachen einer Zahl, so erhält man $16$.
        Wie heisst die Zahl?
  \part Ein Rechteck ist $3$\,cm länger als breit. Sein Umfang beträgt
        $26$\,cm. Wie lang und wie breit ist es?
  \part Ein Kinoticket für Erwachsene kostet CHF 6 mehr als eines für
        Jugendliche. Eine Familie bezahlt für $2$ Erwachsene und $3$
        Jugendliche zusammen CHF 82. Wie viel kostet ein Ticket für
        Jugendliche?
\end{parts}
\begin{solution}
  a) $3x - 5 = 16$, \ $x = 7$. Die Zahl heisst $7$. \
  b) $b$: Breite in cm; $2 \cdot b + 2 \cdot (b + 3) = 26$, \ $4b + 6 = 26$,
  \ $b = 5$. Das Rechteck ist $8$\,cm lang und $5$\,cm breit. \
  c) $j$: Preis für Jugendliche in CHF; $2 \cdot (j + 6) + 3j = 82$,
  \ $5j + 12 = 82$, \ $j = 14$. Ein Ticket für Jugendliche kostet CHF 14
  (Erwachsene CHF 20).
\end{solution}
\raum{6cm}

\question Ein Zug fährt um 7:42 Uhr in Basel ab und legt die $90$\,km bis
Zürich in $54$ Minuten zurück. Unterwegs hält er zweimal.
Berechnen Sie die durchschnittliche Geschwindigkeit in km/h. Streichen Sie
vorher alle Angaben durch, die Sie nicht brauchen.
\begin{solution}
  Überflüssig: Abfahrtszeit und Anzahl Halte. $54$\,min $= 0.9$\,h;
  \ $90 \geteilt 0.9 = 100$. Die durchschnittliche Geschwindigkeit beträgt
  $100$\,km/h.
\end{solution}
\raum{3cm}

\end{questions}

\begin{transfer}
Dieselbe Routine hilft bei Aufgaben in Physik, Chemie und Wirtschaft und
Recht: gegeben, gesucht, Grössen mit Einheiten festlegen, übersetzen, prüfen.
Die Zahlenprobe ist ein allgemeines Werkzeug: Wer unsicher ist, ob eine
Formel stimmt, setzt einfache Zahlen ein und schaut, ob das Ergebnis Sinn
ergibt.
\end{transfer}

\begin{exitfrage}
«In einem Verein gibt es viermal so viele Junioren wie Trainerinnen und
Trainer.» Schreiben Sie dies als Gleichung mit $J$ und $T$ und machen Sie die
Zahlenprobe.
\ifprintanswers\par\medskip\textbf{Lösung:} $J = 4 \cdot T$; Probe mit
$T = 5$: $J = 20$, und es gibt tatsächlich viermal so viele Junioren.
\else\par\vspace{1.5cm}\fi
\end{exitfrage}

% ===========================================================================
\newpage
\werkzeugteil{3}{Von der Mathematik zurück zum Text}
  {ein Rechenergebnis zurück in die Situation übertragen, sinnvoll runden und
   einen Antwortsatz schreiben.}

\begin{einstieg}
Für einen Schulausflug fahren $200$ Personen mit Bussen, die je $45$ Plätze
haben. Der Taschenrechner zeigt: $200 \geteilt 45 = 4.444\ldots$ \\
Wie viele Busse braucht es?
\ifprintanswers\else\par\vspace{1cm}\fi
\end{einstieg}

\begin{lehrperson}[Teil 3, Einstieg]
Die Antworten $4.44$ oder $4$ sind in der Forschung gut dokumentiert: Viele
Schülerinnen und Schüler rechnen korrekt, übertragen das Ergebnis aber nicht
zurück in die Situation. Die Rechnung ist «fertig», sobald eine Zahl
dasteht. Richtig ist: $5$ Busse.
\end{lehrperson}

\begin{merke}
Eine Aufgabe ist erst gelöst, wenn das Ergebnis \textbf{zurück in die
Situation} übertragen ist:
\begin{itemize}
  \item \textbf{Was bedeutet die Zahl?} Welche Einheit hat sie?
  \item \textbf{Sinnvoll runden:} Die Situation entscheidet, ob auf- oder
        abgerundet wird.
  \item \textbf{Plausibel?} Ist das Ergebnis möglich und von der
        Grössenordnung her vernünftig? Eine negative Anzahl Personen oder ein
        Schulzimmer mit $18$\,m$^3$ Volumen deutet auf einen Fehler hin.
  \item \textbf{Antwortsatz:} Die Frage in einem ganzen Satz beantworten,
        mit Einheit.
\end{itemize}
\end{merke}

\begin{beispiel}[1: Aufrunden oder abrunden?]
\begin{itemize}
  \item $200$ Personen, Busse mit $45$ Plätzen: $200 \geteilt 45 \approx 4.4$.
        \textbf{Aufrunden:} Es braucht $5$ Busse, sonst bleiben Personen
        zurück.
  \item Ein Seil von $10$\,m wird in Stücke von $3$\,m geschnitten:
        $10 \geteilt 3 \approx 3.3$. \textbf{Abrunden:} Es entstehen $3$ ganze
        Stücke; der Rest von $1$\,m ist zu kurz.
\end{itemize}
Dieselbe Rechnung, eine andere Situation, eine andere Rundung.
\end{beispiel}

\begin{beispiel}[2: Ein Ergebnis vollständig deuten]
Handyabo A kostet CHF 20 pro Monat plus CHF 0.10 pro Minute, Abo B pauschal
CHF 35. Aus $20 + 0.1x = 35$ folgt $x = 150$.

\falsch\quad «$150$»

\richtig\quad «Bei genau $150$ Minuten kosten beide Abos gleich viel. Wer mehr
als $150$ Minuten pro Monat telefoniert, fährt mit Abo B günstiger, wer
weniger telefoniert, mit Abo A.»
\end{beispiel}

\begin{beispiel}[3: Terme in Worten lesen]
\begin{center}
\renewcommand{\arraystretch}{1.4}
\begin{tabular}{ll}
  $3 \cdot (x + 4)$ & das Dreifache der Summe von $x$ und $4$ \\
  $3x + 4$          & die Summe aus dem Dreifachen von $x$ und $4$ \\
  $\dfrac{a + b}{2}$ & die Hälfte der Summe von $a$ und $b$ (der Mittelwert)
\end{tabular}
\end{center}
Wer einen Term in Worten lesen kann, erkennt seine Struktur: Was wird
\emph{zuletzt} gerechnet? Danach richtet sich der Name (Summe, Produkt, \ldots).
\end{beispiel}

\begin{questions}

\question Lösen Sie die Aufgaben und schreiben Sie jeweils einen Antwortsatz.
\begin{parts}
  \part Eine Dose Farbe reicht für $12$\,m$^2$. Eine Wand ist $50$\,m$^2$
        gross. Wie viele Dosen braucht es?
  \part Ein Lift darf höchstens $630$\,kg tragen. Wie viele Kisten zu je
        $75$\,kg dürfen gleichzeitig transportiert werden?
  \part Beim Lösen einer Textaufgabe zur Anzahl Personen in einem Raum
        erhält Mia $x = -3$. Was bedeutet das?
\end{parts}
\begin{solution}
  a) $50 \geteilt 12 \approx 4.2$; es braucht $5$ Dosen (aufrunden). \
  b) $630 \geteilt 75 = 8.4$; es dürfen höchstens $8$ Kisten transportiert
  werden (abrunden). \
  c) Eine negative Anzahl Personen ist nicht möglich. Entweder hat Mia einen
  Fehler gemacht (Gleichung oder Rechnung prüfen), oder die Situation hat keine
  Lösung.
\end{solution}
\raum{4cm}

\question Beschreiben Sie die Terme in Worten.
\begin{multicols}{2}
\begin{parts}
  \part $a + \dfrac{b}{2}$
  \part $2 \cdot (x - 1)$
  \part $2x - 1$
  \part $(a + b)^2$
\end{parts}
\end{multicols}
\begin{solution}
  a) die Summe aus $a$ und der Hälfte von $b$ \
  b) das Doppelte der Differenz von $x$ und $1$ \
  c) die Differenz aus dem Doppelten von $x$ und $1$ \
  d) das Quadrat der Summe von $a$ und $b$
\end{solution}
\raum{3cm}

\question Erfinden Sie eine Situation aus dem Alltag, die zur Gleichung
$5x + 20 = 95$ passt. Lösen Sie die Gleichung und schreiben Sie einen
Antwortsatz zu Ihrer Situation.
\begin{solution}
  $x = 15$. Mögliche Situation: Ein Fitnessabo kostet eine einmalige Gebühr
  von CHF 20 und CHF 5 pro Besuch. Jemand hat insgesamt CHF 95 bezahlt. Wie oft
  war die Person im Fitness? Antwort: Die Person war $15$-mal im Fitness.
\end{solution}
\raum{4cm}

\end{questions}

\begin{lehrperson}[Übung 3]
Übung 3 ist die Umkehrung der üblichen Textaufgabe und für viele
ungewohnt. Sie zeigt sehr direkt, ob jemand versteht, was die Teile einer
Gleichung bedeuten: Die $5$ muss ein Preis pro Stück oder eine ähnliche
Rate sein, die $20$ ein fester Betrag. Gute Beispiele aus der Klasse eignen
sich als Aufwärmaufgaben für spätere Lektionen.
\end{lehrperson}

\begin{transfer}
Auch in den Naturwissenschaften ist eine Rechnung erst fertig, wenn das
Ergebnis gedeutet ist: Eine Geschwindigkeit von $3000$\,km/h für ein Velo
oder eine Masse von $-2$\,kg zeigt sofort, dass etwas nicht stimmt. Und in
jedem Fach gilt: Wer die Frage mit einem ganzen Satz beantwortet, merkt eher,
ob er sie überhaupt beantwortet hat.
\end{transfer}

\begin{exitfrage}
$95$ Schülerinnen und Schüler sollen in Gruppen zu höchstens $6$ Personen
eingeteilt werden. Wie viele Gruppen braucht es mindestens? Antworten Sie mit
einem ganzen Satz.
\ifprintanswers\par\medskip\textbf{Lösung:} $95 \geteilt 6 \approx 15.8$. Es
braucht mindestens $16$ Gruppen.
\else\par\vspace{1.5cm}\fi
\end{exitfrage}

\end{document}
